\paragraph{Passing the detecting count through the auxiliary tower.}
Let $\ell=\operatorname{char}k\ne2$, and retain the original finite
paired modules $W_m$. Write $\Gamma=G\rtimes\langle c\rangle$,
$c^2=1$, and let $\beta$ be their original action, with
\[
 \beta(g)Z=\chi(g)R(g)ZR(g)^{-1},\qquad
 \beta(c)Z=-\chi(c)Z^T.
\]
For an open normal subgroup $H$ of a profinite group, write
$P_\ell(H)$ for the kernel of its maximal pro-$\ell$ quotient.
The original subgroup defined using ambient open normal quotients is
the same: taking the finite normal core of a kernel of an
$\ell$-group quotient of $H$ preserves an $\ell$-group quotient.

Retain the original tower $p_N:\Gamma_N\twoheadrightarrow\Gamma$
and its compatible involutions. The action upstairs is the pullback
$\beta\circ p_N$. Let $H$ be the original open normal subgroup of
$\Gamma$, with $H\subset G\cap\ker\chi$, and retain
$\rho(P_\ell(H))=1$. Put $H_N=p_N^{-1}(H)$.
The actual tower conditions used here are
\[
 \begin{gathered}
 p_N(\tau_{N,j})=1,\\
 \ker p_N\subseteq
 \overline{\langle\!\langle\tau_{N,j}\rangle\!\rangle P_\ell(H_N)},\\
 p_N(\varphi_{N,j})\sigma_j^{-1}\longrightarrow1
                  \text{ in }\Gamma/P_\ell(H).
 \end{gathered}
\]
where $\sigma_j$ are the fixed detecting elements in the cyclotomic
kernel. Assume the original finiteness condition on continuous
homomorphisms from open subgroups to $\mathbb F_\ell$.

\begin{proposition}[The unchanged detecting constant]
\label{prop:selmer-tower-detection}
Let $D_m$ be the detecting kernel among polarized paired base
cocycles, and $B_m$ its polarized boundaries. If
$\#(D_m/B_m)\leq q^C$ for all $m$, then, for each $m$ and all
sufficiently large $N$, the set $S_{N,m}$ of continuous cocycles
satisfying
\[
 \forall j\ \exists Y_j:\quad
 y(\varphi_{N,j})=\delta Y_j(\varphi_{N,j}),\qquad
 y(\tau_{N,j})=\delta Y_j(\tau_{N,j})
\]
is finite and obeys
\[
        \#S_{N,m}\leq\#B^1(\Gamma_N,W_m)\,q^C.
\]
The same $C$ supplied by integral detection and
Lemma~\ref{shared-complete-solution-count} therefore works throughout
the tower. Only the level threshold depends on $m$.
\end{proposition}
\begin{proof}
At fixed $m$, the image of $H$ under $\rho$ is an $\ell$-group.
Passing to $R$ adds a subgroup of
$1+\sum_i\varepsilon_iM_n(A/(t^m))$, an elementary abelian
$\ell$-group. Since $\chi|_H=1$, the image of $H$, and of $H_N$,
under the paired action is also an $\ell$-group.

For a cocycle $y$, the affine action $v\mapsto y(g)+\beta(g)v$
on $H_N$ extends this image by a subgroup of the additive
$\ell$-group $W_m$. Its open normal kernel contains $P_\ell(H_N)$.
If $y(\tau_{N,j})=0$, it also contains every $\tau_{N,j}$:
both the representation and twist are trivial there. The tower
kernel condition therefore makes $y$ descend continuously to the
base. Continuity follows from the compact-to-Hausdorff quotient map.

There are only finitely many base cocycles at fixed $m$: restrict to
an open kernel of the finite action and then choose finitely many
coset values. The original finiteness condition controls the
homomorphisms on that kernel. For each such cocycle its affine
action on $H$ has $\ell$-group image. Intersect the finite normal
cores of its kernel with the representation and character kernels.
This gives an ambient open normal subgroup containing $P_\ell(H)$.
Tower convergence consequently implies, for all large $N$ and
simultaneously for every base cocycle $z$ and every $j$,
\[
 R(p_N\varphi_{N,j})=R(\sigma_j),\qquad
 \chi(p_N\varphi_{N,j})=1,\qquad
 z(p_N\varphi_{N,j})=z(\sigma_j).
\]

Every $y\in S_{N,m}$ vanishes at $\tau_{N,j}$ and hence descends.
The identity $\beta(c_N)y(c_N)=-y(c_N)$ gives
\[
 v=-y(c_N)/2,\qquad \delta v(c_N)=y(c_N).
\]
Thus $y^\circ=y-\delta v$ is zero at $c_N$, has local boundary
primitives $Y_j-v$, and descends to a polarized base cocycle
$z$. At a detecting element, the eventual equalities give
$z(\sigma_j)=R(\sigma_j)Y_jR(\sigma_j)^{-1}-Y_j$.
Cyclicity of trace makes
$[\eta]\operatorname{tr}(z(\sigma_j)R(\sigma_j))=0$.

This induces an injection
\[
       S_{N,m}/B^1(\Gamma_N,W_m)\hookrightarrow D_m/B_m.
\]
Indeed, normalizing a boundary $\delta w$ replaces its primitive
by $(w+\beta(c_N)w)/2$, which is fixed by $c_N$. Conversely a
polarized primitive for $z$ gives a global boundary agreeing with
$y^\circ$ on the index-two subgroup and at $c_N$, hence everywhere.
The right quotient is bounded by $q^C$ and the boundary group is
finite. This proves both assertions, with no additional factor
\cite{alpo-tower-detection}.
\end{proof}

The argument uses all modular cocycles and retains the original
kernel and convergence conditions. It does not assume that finite
cocycles lift, nor construct the arithmetic tower or its points.
